\documentclass[11pt,a4paper]{article}
\usepackage[margin=27mm]{geometry}
\usepackage{amsmath,amssymb,amsthm,hyperref}
\newtheorem{lemma}{Lemma}
\newcommand{\Sz}{\operatorname{Sz}}
\newcommand{\F}{\mathbb F}
\newcommand{\gen}[1]{\langle#1\rangle}
\title{Preliminary statements and their precise scope}
\author{Human Compression project}\date{}
\begin{document}\maketitle
\begin{abstract}
We record the complete preliminary dependency boundary for the fixed 1999
rank-three Suzuki article. Several literal arithmetic assertions need
correction: integer divisibility is insufficient for subfield divisibility,
the residue class controlling the factor divisible by five is misstated,
and a prime extension must remain in the same torus. We give small exact
witnesses and the corrected general statements. The standard group structure
and the explicitly cited linear-diagram theorem remain attributed inputs.
\end{abstract}

\section{Definitions and quoted structural facts}
Throughout, the source is D.~Leemans, \emph{The rank 3 geometries of the
simple Suzuki groups $\Sz(q)$}, Note di Matematica 19(1) (1999), 43--63,
in the fixed twenty-one-page edition. Write
\[
 q=2^{2e+1},\qquad \theta=2^{e+1},\qquad
 \alpha_q=q+\theta+1,\quad \beta_q=q-\theta+1.
\]
A simple Suzuki subfield group has size parameter at least eight.
The source uses $a\mid_M b$ to mean that $b/a$ is prime, and
$a\mid_P b$ to mean that $a$ is a proper divisor of $b$ (printed p.~47).
These symbols must not be replaced by ordinary divisibility in a source join.

Source Lemma~2.1 states that a rank-three coset pregeometry is a geometry.
Indeed an incident pair of cosets has a common representative $g$, and
adding $gG_k$ of the missing type gives a chamber. A singleton similarly
extends using any representative; the empty flag extends to the base
chamber. Thus every flag extends to a chamber. No flag-transitivity
assertion follows from this argument.

The classical Suzuki maximal-subgroup list in source Table~1 is an
explicitly attributed premise from L\"uneburg and Suzuki. Source
Lemmas~3.1 and~3.2 are explicitly quoted from the earlier rank-two paper.
Their required content is independently derived in the accepted L115
preliminary note: the maximal subgroups of $B_q,D_q,N_+(q),N_-(q)$,
and pairwise coprimality of $q-1,\alpha_q,\beta_q$.
The resulting bounds on pairwise intersection orders, respectively
$2,2,4$, follow from divisibility of subgroup orders.
No additional finite census is used here.

Source Lemma~3.8 is expressly stated to follow from the main theorem
of M.~Aschbacher, \emph{Flag structures on Tits geometries}, Geometriae
Dedicata 14 (1983), 21--31. We preserve this as a cited premise in its
source form: a rank-three geometry with a linear diagram is flag
transitive. This note does not reconstruct the cited theorem. In particular,
no arbitrary intersection pattern is declared to have such a diagram
without the relevant residue and diagram hypotheses.

\section{Field degrees and the arithmetic statements}
The classical list gives one conjugacy class of maximal simple subfield
groups for each prime $p\mid n$, where $q=2^n$, provided $n/p\ge3$.
If $n$ is prime there are none. If $n$ is composite all these quotients
are at least three. Consequently there is exactly one class if and only
if $n=p^k$, $k\ge2$. This proves source Lemma~3.3.

The repunit argument in the accompanying normalizer-restriction note
proves the implication in Lemma~3.4: if $\Sz(s)$ is maximal in $\Sz(q)$
and $(q-1)/(s-1)$ is prime, then $q=s^p$, $s=2^{p^a}$ for a prime
$p$ and a positive integer $a$. Therefore, if there are at least two
maximal subfield classes, $s-1$ is \emph{not a maximal proper divisor}
of $q-1$. This is the consequence needed in the source.

\paragraph{Lemma 3.5: the printed word ``prime''.}
Printed p.~48 instead says that $s-1$ is not a prime divisor of $q-1$.
That literal assertion is false. Take $q=2^{15}=32768$ and $s=8$.
The subfield group $\Sz(8)$ is maximal, of field index five, and there
is also the different maximal class $\Sz(32)$, of field index three.
Yet $s-1=7$ is prime and
\[
 32767=7\cdot31\cdot151.
\]
Its quotient $32767/7=4681$ is composite, in agreement with the corrected
maximal-divisor assertion. This is a wording correction to an ancillary
lemma, not an additional principal-refutation count.

\paragraph{Lemma 3.6: divisibility and the residue variable.}
Part (1) is $5\mid q^2+1$, not $s^2+1\mid q^2+1$.
It follows at once from $q^2=2^{4e+2}\equiv-1\pmod5$.
For part (2), the native statement assumes merely $s\ge8$ divides $q$.
This is insufficient. At $s=8,q=32$ it holds as integer divisibility,
but
\[
 (\alpha_8,\beta_8)=(13,5),\qquad
 (\alpha_{32},\beta_{32})=(41,25),
\]
and $13$ divides neither $41$ nor $25$. The corrected hypothesis is
$s=2^d$, $d\mid 2e+1$, equivalently $q=s^r$ for an odd integer $r$.
Under that hypothesis the compatible-torus divisibilities, with the
possible interchange of the two signs, are proved in the accepted L115
preliminary note. Only that version is imported in our proofs.

Part (3) prints the residue class of $2e+1$ modulo four. The correct
variable is $e$:
\[
 5\mid\alpha_q\ \Longleftrightarrow\ e\equiv0,3\pmod4,
 \qquad
 5\mid\beta_q\ \Longleftrightarrow\ e\equiv1,2\pmod4.
\]
To prove this for all $e$, observe that both $2^{2e+1}$ and $2^{e+1}$
are periodic modulo five with period four in $e$. For $e=0,1,2,3$,
the respective residues of $(\alpha_q,\beta_q)$ are
$(0,1),(3,0),(1,0),(0,3)$. The printed formula already fails at
$q=8$, where $2e+1=3$ but $\alpha_q=13$ is not divisible by five.

\begin{lemma}
For odd positive integers $m,n$,
$\gcd(2^{2m}+1,2^{2n}+1)=2^{2\gcd(m,n)}+1$.
\end{lemma}
\begin{proof}
Put $x=4$ and $d=\gcd(m,n)$. The Euclidean algorithm for powers
gives $\gcd(x^{2m}-1,x^{2n}-1)=x^{2d}-1$.
The common divisor $D=\gcd(x^m+1,x^n+1)$ therefore divides
$(x^d-1)(x^d+1)$. It is coprime to $x^d-1$: the latter divides
$x^m-1$, and an odd common divisor of $x^m-1,x^m+1$ is one.
Hence $D\mid x^d+1$. Conversely $m/d,n/d$ are odd, so
$x^d+1$ divides both $x^m+1$ and $x^n+1$. This proves the equality
and source Lemma~3.7.
\end{proof}

\section{Unique extensions and proper overgroups}
For source Lemma~3.9(i), let $H=C_d\rtimes C_2$, $d>1$ dividing
$q^2+1$. The cyclic odd subgroup determines one ambient nonsplit torus
$T$ and its normalizer $N=T\rtimes C_4$. Every group of the form
$C_d\rtimes C_4$ containing $H$ normalizes its characteristic odd
subgroup and hence lies in $N$. An involution of $N$ is contained in
exactly one cyclic subgroup of order four: that cyclic subgroup is its
centralizer in $N$, since the involution inverts $T$. Adjoining this
$C_4$ to the unique order-$d$ subgroup of $T$ gives the unique extension.
This proves (i), including actual existence.

For (ii), write $|T|=a_\epsilon(q)$. If $H=C_d\rtimes C_4$ and
$p\mid |T|/d$ is prime, the unique subgroup of $T$ of order $dp$,
together with the given $C_4$, supplies a unique $C_{dp}\rtimes C_4$
containing $H$. Every possible extension must have this torus and
this cyclic odd subgroup, proving uniqueness.

\paragraph{Lemma 3.9(ii): the required compatible factor.}
The native condition is instead $p\mid(q^2+1)/d$, which can choose
a prime from the opposite, coprime nonsplit factor. At $q=8,d=5,p=13$,
the subgroup $C_5\rtimes C_4$ exists, but $C_{65}\rtimes C_4$
does not: it would contain an element of order65, whereas the odd
elements of $\Sz(8)$ lie in cyclic tori of orders $5,7,13$.
The source condition therefore needs the compatible-factor restriction
just proved. The retained torus partition is the only group-theoretic
premise in this small obstruction; no enumeration is required.

The full split-torus overgroup argument in the normalizer-restriction
note gives the precise version of source Lemma~3.10: the \emph{strict
proper} overgroups of $C_{q-1}$ are its dihedral normalizer, the two
point stabilizers, and their two $E_q\rtimes C_{q-1}$ subgroups.
The torus itself and the whole ambient group must be added if neither
strictness nor properness is imposed. We state this qualification
explicitly rather than assigning an unintended meaning to ``containing''.

Finally let $S=\Sz(s)$ be maximal in $G=\Sz(q)$ and let
$H=D_{2(s-1)}$ be its full split normalizer. Every proper overgroup of
$H$ is either a dihedral subgroup $D_{2u}$ of its ambient split normalizer,
with $s-1\mid u\mid q-1$, or a simple subfield group $\Sz(t)$.
Point stabilizers are excluded by the inversion argument, and nonsplit
normalizers by coprimality. In the subfield case the rotation subgroup
is split, so $s-1\mid t-1$. Writing $s=2^d,t=2^f$ gives $d\mid f$;
this follows from the Euclidean identity
$\gcd(2^d-1,2^f-1)=2^{\gcd(d,f)}-1$.
Maximality of the field degree $d$ in the ambient degree forces $f=d$
for a proper subfield. The accepted uniqueness of the subfield containing
this embedded full split normalizer then gives $\Sz(t)=S$.
This proves source Lemma~3.11 for proper overgroups. The ambient group
itself remains excluded by that explicit qualification.

\section{Attribution and remaining work}
The new source joins above distinguish literal errors, intended corrected
statements, and earlier quoted inputs. They do not replace the primary
classification counterexamples already accepted for L114, and add no
principal-refutation or ordinary-recovery count. The arithmetic checks
attached to this note only verify the displayed small integers; the general
proofs and imported subgroup premises require distinct mathematical review.
The source's general rank-three thinness assertion under its literal
algebraic local-action convention and the final source-case joins remain
separate obligations. No Lean verification or compression measurement is
claimed.
\end{document}
